\documentclass[10pt]{article}
\usepackage{icmlstyle}
\usepackage{booktabs}
\usepackage{amsmath}
\usepackage{amssymb}
\usepackage[T1]{fontenc}
\usepackage[utf8]{inputenc}
\usepackage{url}
\usepackage{xcolor}

\begin{document}

\icmltitle{Five Levers, One Uninterpretable Result, and the Instrument\\[0.2ex]
Required to Tell Them Apart}

\icmlauthor{OPHIS Autonomous Research Campaign}%
{Single-node H200, frozen data pipeline, seed 42}

\begin{icmlabstract}
We report six interventions measured on a fixed-\emph{time} language-model training
benchmark: a $\sim$50.3M-parameter GPT trained for 300 charged seconds on one H200 with
the data pipeline frozen. Three lower validation bits-per-byte and three do not. The
largest single effect is HALVING the tokens processed per optimizer step, which lowers
$\mathrm{val\_bpb}$ by $0.007953$ and produces our best model at $0.984017$, down from
$0.990345$. We report it as a configuration result and explicitly NOT as a statement about
batch size: halving the batch doubles the step count, and at least six quantities in the
trainer are indexed by step rather than by elapsed time, so six things moved at once. Of
the remainder, halving a sliding-window attention span pays twice over two rungs, and
stacking three separately confirmed levers realises about $89\%$ of their additive sum.

Our central methodological claim is that none of these effects is visible without a
paired, device-counterbalanced design: the 70 byte-identical control runs in this campaign
span a range of $0.034609$, roughly an order of magnitude larger than the largest effect
we report. We describe the design, the noise model motivating it, and---at length, because
we believe it is the more transferable contribution---the instrument defects that had to
be found and corrected before any number below could be trusted. Six of the seven defects
repaired in this block were in our own apparatus rather than in the model.
\end{icmlabstract}

\begin{multicols}{2}

\section{Introduction}

The benchmark fixes \emph{time}, not tokens. A change that improves the model per step but
costs throughput can be a net loss, and a change that buys steps by removing computation
can be a net win even if it degrades the model per step. Throughput is therefore part of
quality rather than a nuisance to be regressed away. We report raw $\mathrm{val\_bpb}$ as
the verdict everywhere and use step count only to explain \emph{why} a result came out as
it did.

The campaign began with no experimental priors; every inherited belief was deleted
deliberately, on the grounds that most were wrong or unfalsifiable. What we had was a
literature corpus and a set of properties of the apparatus. What we lacked---and had to
build first---was a measuring instrument able to resolve the effects we sought.

\paragraph{Contributions.}
(i)~Three confirmed levers and three confirmed non-levers, each measured under a
counterbalanced design with activation independently verified (\S\ref{sec:design},
\S\ref{sec:results}).
(ii)~A combination experiment showing the levers are complementary but partly redundant,
with the overlap channel identified (\S\ref{sec:stack}).
(iii)~An account of how the measurement apparatus failed, and what that implies for
autonomous experimentation (\S\ref{sec:instrument}).

\section{Setup}

A GPT of $\sim$50.3M parameters: vocabulary 8192, depth 8, $n_{\text{embd}}$ 512, MLP
ratio 4, head dimension 128. Attention follows a repeating \texttt{SSSL} window
pattern---three short-window layers and one long-window layer per block, final layer
forced to full context. Value embeddings attach to every other layer; QK-norm is applied.
The optimizer is Muon on 2D matrices with AdamW elsewhere. Training runs under
\texttt{torch.compile} in bf16 for 300 charged seconds: on this host roughly 1000 steps at
524288 tokens per step, completing about two epochs of a fixed dataset.

The tokenizer, loader, packing, evaluator and BPB accounting are frozen and byte-identical
to the upstream reference. Only the training file is edited. The host is shared; other
tenants occupy most GPUs most of the time, and co-tenancy on a GPU invalidates a run.

\section{The measurement problem}

The 70 byte-identical control runs span $1.025517$ at worst to $0.990908$ at best---a
range of $0.034609$. Every effect we report is smaller than $0.004408$. A single unpaired
run therefore cannot see any of them.

The dominant nuisance is host CPU contention, which moves step count run-to-run through
the frozen packing loop. Two structural facts follow.

\paragraph{Concurrency, not proximity.} Runs must overlap in wall-clock time to share
contention. Runs launched back-to-back on one GPU are sequential and the load drifts
between them; grouping those as a ``wave'' reports a sequential spread under a within-wave
label, flattering the instrument.

\paragraph{The offset belongs to the GPU.} Our four devices have systematically different
control means, with a fastest-to-slowest spread comparable to the effects being chased. An
earlier version of this campaign attributed that offset to the CPU core block; correcting
which factor owns it changed verdicts.

\section{Design: the counterbalanced quad}
\label{sec:design}

Each experiment is two four-wide waves launched concurrently:
\begin{center}\small
\begin{tabular}{ll}
wave A: & [\,treat, ctrl, ctrl, treat\,]\\
wave B: & [\,ctrl, treat, treat, ctrl\,]
\end{tabular}
\end{center}
Across the pair the treatment occupies every device exactly once, as does the control.
Deltas are formed \emph{within} each device, removing that device's fixed offset inside
each difference rather than relying on cancellation in the mean.

Two consequences matter. First, $n=4$ pairings, so every $t$ carries 3 degrees of freedom
and should be read as evidence of \emph{direction}, not as a precise magnitude. Second,
when the design breaks---a device set shifts between waves, or an arm crosses an
operating-point boundary---the analysis \textbf{refuses a verdict} rather than averaging
across a broken swap. This occurred three times in this block, costing three extra waves.

\paragraph{Activation.} Every hypothesis declares a diagnostic the run must emit, checked
\emph{before} the outcome. A run whose diagnostic is absent, or which fails its rule, is
\textbf{inconclusive} about the mechanism---never a negative. Otherwise a broken
intervention and a genuine null are indistinguishable and good ideas are retired for free.

The instrument resolution, twice the pooled within-device control standard deviation
divided by $\sqrt{2}$, is $0.000282$ (pooled sd $0.000199$). Effects below it are not
claimed.

\section{Results}
\label{sec:results}

All comparisons are same-device paired within a single counterbalanced quad pair, seed 42.

\begin{center}\small
\begin{tabular}{@{}lrr@{}}
\toprule
intervention & mean $\Delta$ & verdict\\
\midrule
\texttt{tbs} 19$\to$18 & $-0.007953$ & better\textsuperscript{\dag}\\
\texttt{swdiv} 4$\to$8 & $-0.003360$ & better\\
\texttt{swdiv} 2$\to$4 & $-0.002298$ & better\\
\texttt{ve} 2$\to$1 & $-0.001487$ & better\\
\texttt{precond} & $-0.001147$ & better\\
\texttt{qk\_suppress} 0.1 & $+0.001226$ & \emph{inconclusive}\\
\texttt{unet} & $+0.004123$ & worse\\
\texttt{tbs} 19$\to$20 & $+0.022545$ & worse\\
\texttt{mtp} 0.1 & $+0.079484$ & worse\textsuperscript{\ddag}\\
\bottomrule
\end{tabular}\\[0.4ex]
{\footnotesize \textsuperscript{\dag}confounded; see \S\ref{sec:sixlegs}.\quad\textsuperscript{\ddag}three replicates; see \S\ref{sec:mtp}.}
\end{center}

\subsection{Halving the short attention span}

At sequence length 2048 the \texttt{SSSL} pattern gives six short-window and two
long-window layers. Halving the short window from 1024 to 512 takes per-token FLOPs from
$239.078$ to $220.204$: a removal of $18.874$, or $7.9\%$ of the total---\emph{not} the
$26\%$ that attention as a whole represents. The two long-window layers hold about $40\%$
of the attention term and are untouched.

This distinction was not academic. The hypothesis was originally justified with the $26\%$
figure, overstating the throughput lever by $\sim$3.3$\times$. The literature review caught
the error before the experiment ran; the arithmetic was corrected and the prediction
threshold was \emph{deliberately not lowered} to compensate. The corrected, harder
prediction was met anyway.

Step counts rose in every cell and quality held. This is the opposite outcome to an
earlier arm reducing Newton--Schulz iterations, worth $\sim$0.2\% of step time, which
bought no steps and came back worse: $7.9\%$ is enough to move the critical path and
$0.2\%$ is not.

\paragraph{A correction we owe the reader.} We first reported this effect as $-0.002463$
at $t=-20.3$. An independent audit found that estimate stitched pairs from two experiments
run 1.4 hours apart. The clean single-experiment quad gives $-0.002298$ (sd $0.000142$,
$t=-32.4$). The stitched figure overstated the effect by about $7\%$; the corrected value
is used throughout.

\subsection{Value embeddings on every layer}

\texttt{ve} is a \emph{stride}, so \texttt{ve=1} is denser: tables go from four to eight
and parameters from 50.3M to 67.1M. This 33\% larger model ran slower and completed
\emph{fewer} steps than its control in every cell, and won regardless---quality gained per
step more than repays throughput lost.

This contradicts the prior we registered before running it. Three corpus claims opposed
it: all-layer injection measured worse than single-layer, dense per-layer expansion
saturating, and evidence that such memories' benefit is the gated pathway rather than the
learned table content. All three carried honestly low transfer scores (image
autoregression; tokenizer-transfer fine-tuning; a 2.5B mixture-of-experts). They retain
their opposing stance; what we now know is that they do not carry to 50M parameters at 300
seconds.

\paragraph{One void cell.} On the slowest device the added parameters cost enough
throughput to finish at $501.2$M tokens and epoch $1.0$ against its control's $2.0$---a
different operating point, so the pair is excluded. The exclusion rule is structurally
biased against parameter-adding treatments; here it cut \emph{against} the result, since
including the cell would have strengthened it. This leaves \texttt{ve} with three pairings
and 2 degrees of freedom, the weakest of our three wins.

\subsection{Second-moment reordering}

Moving the NorMuon second-moment computation before the polar iteration, with unit-RMS
momentum normalisation, gives $-0.001147$ across four devices. The activation observable
separates the arms by roughly two orders of magnitude, because the second moment is taken
over raw momentum rather than an orthogonalized update.

\subsection{Two clean negatives}

\textbf{Zero-initialised U-net skips} raised $\mathrm{val\_bpb}$ by $0.004123$ with a
four-device standard deviation of only $0.000078$---the tightest measurement in the
campaign. That precision follows from the intervention: the skips add one elementwise
addition per lower-half layer and no matrix multiply, so they cannot buy or cost steps,
and the step counts confirm they did not. The effect is pure quality, free of the
throughput entanglement that explains most variance elsewhere. The gate parameters grew
from exactly zero to $\sim$0.2, so the model \emph{used} the path and was worse for it.
This intervention was the independent top pick of two separate proposers in the same
council round; their agreement carried no evidential weight whatever.

\textbf{Doubling tokens per step} raised $\mathrm{val\_bpb}$ by $0.022545$, roughly eighty
times the resolution and by far the largest effect measured. Step counts halved as
expected. At fixed time this model sits far \emph{above} its critical batch size, so
halving the optimizer update count costs far more than reduced gradient noise returns. The
step law cannot explain this away: it was fitted at 524288 tokens per step and this arm
moves that very quantity.

\subsection{A mechanism that worked and lost anyway}
\label{sec:mtp}

Multi-token prediction is the largest effect this campaign has measured and it is a
\emph{negative}: $+0.079484$, roughly thirty times the largest slot offset ever observed on
this host, reproduced across three independent pairs whose deltas agree to $6\times10^{-4}$
($+0.079681$, $+0.079288$, $+0.079131$).

It is worth reporting because the mechanism plainly \emph{engaged}. The auxiliary head's
loss fell by $1.967$ nats over the run against an activation threshold of $0.1$, so this is
a valid negative and not a non-activation: the $t+2$ objective learned, and the model was
still worse. The cost is entirely throughput. The extra $[B,T,V]$ backward pass through the
vocabulary head ran at $834$\,ms per step against the control's $300$\,ms, delivering $371$
optimizer steps against $1011$, and reserving $107.5$\,GB against $45.1$\,GB. At a fixed
$300$-second budget the treatment never reached the second epoch its control completed.

The result generalises less far than its size suggests. It says nothing about MTP at larger
vocabularies, longer budgets, or with a factorised head; it says that at this operating
point a mechanism which doubles per-step cost cannot be rescued by being right. The
practical rule we now apply is that any intervention adding a full-width logits pass is
screened on step count \emph{before} it is screened on \texttt{val\_bpb}.

This result also broke the reporting tool. Our verdict code voids any comparison whose
treatment and control finish at different epochs, on the reasonable ground that a run which
crossed the two-epoch boundary is a different operating point. Here \emph{every} arm showed
the gap, in the same direction, because the treatment was too slow to reach the boundary at
all --- and the tool printed no verdict for the largest effect on record. Direction
consistency alone is not evidence of cause, so the rule now additionally requires a large
throughput change and an effect the instrument can resolve; a sub-band delta with a $1\%$
step wobble straddling the boundary is still excluded.

\subsection{One honest inconclusive}

\texttt{qk\_suppress} produced a textbook-looking counterbalanced negative, $+0.001226$ at
$t=+15.8$. We do \emph{not} report it as a negative. Its declared diagnostic reads exactly
$1.0$ in every treatment \emph{and} every control, because QK-norm pins query RMS to unit
norm by construction. It could never have shown engagement, so the run cannot distinguish
``the mechanism works badly'' from ``the edit never bound and something else moved''. The
hypothesis's own pre-registered falsifier named this case and we apply it against
ourselves. The axis is untested, not refuted.

\section{Combining the levers}
\label{sec:stack}

The three winners were run together against the \emph{unchanged} control:
\[
\Delta = -0.004408,\quad \mathrm{sd}=0.000152,\quad \mathrm{sem}=0.000076
\]
over four devices, all eight arms at epoch $2.0$, with all three components independently
confirmed active in the same records. The best model is $0.986956$.

The naive sum of the three separately measured effects is $-0.004932$. The combination
realises about $89\%$ of it, a shortfall of $0.000524$.

\paragraph{Where the overlap is.} \texttt{ve} and \texttt{swdiv} both act partly through
step count, in opposite directions: \texttt{ve} costs throughput by adding parameters,
\texttt{swdiv} buys it by removing FLOPs. Stacked, they nearly cancel in that
channel---on the slowest device the combined arm ran 991 steps against its control's 992,
where \texttt{ve} alone managed 956 and voided. Some of the two levers' separate gains
were the same gain counted twice.

\paragraph{Why this reading is licensed.} The interaction was predicted \emph{before} the
run and the reading order fixed in advance: check epoch on every arm first, then look at
$\mathrm{val\_bpb}$. Had the combined arm voided on the slow device, a shortfall would
have been an operating-point artefact saying nothing about redundancy. It did not void,
and the step count shows why. After the fact the two explanations are indistinguishable,
which is the entire reason the order was pre-committed.

\paragraph{On the statistics.} Our reported $t=-58$ rides a 3-df variance estimate about
half the control-derived expectation; a chi-square check gives $p\approx0.14$, so this is
not evidence of genuinely lower treatment noise, and a fairer $t$ is $\approx-29$. The win
is unambiguous either way. We initially justified the shortfall as ``2.4$\times$
resolution'', which repeats a borrowed-constant error our own lesson log condemns: the
additive sum is itself three noisy estimates and its error must be propagated. Doing so
properly gives $\approx$3.3 standard errors at Welch $\mathrm{df}\approx8.9$, so
sub-additive redundancy stands---but by correct arithmetic, not the argument we first
gave.


\paragraph{Two caveats an audit raised after first submission.} First, \texttt{precond}'s
$t=-12.3$ is a same-device paired statistic computed \emph{across} waves; the raw
within-wave paired $t$ is $\approx-2.0$. Both share the mean $-0.001147$, and the
same-device pairing is defensible because it removes the device term inside each
difference rather than leaving it to cancel---but it relies on the device model being
stable, so it is model-adjusted rather than raw, and we state it as such. Second, the
configuration string \texttt{precond="pre"} names \emph{two} different implementations in
our records, the pre-repair version whose second-moment buffer collapsed onto its clamp
floor and the repaired one; only the latter supports the number above, and our analysis
keys arms by generated-source hash precisely so the two are never pooled.

\section{The largest result is uninterpretable}
\label{sec:sixlegs}

Halving tokens per step from $524288$ to $262144$ lowered $\mathrm{val\_bpb}$ by
$0.007953$ (four devices, all at epoch $2.0$), roughly twice the three-lever stack and by
far the largest effect we have measured. We can state the configuration result. We cannot
state the mechanism, and the distinction is the point of this section.

Halving the batch doubles the step count, from about $990$ to about $2000$. The trainer's
learning-rate and weight-decay schedules are indexed by \emph{progress} -- elapsed time
over budget -- but at least six other quantities are indexed by \emph{step}, so all of
them moved simultaneously:

\begin{enumerate}\itemsep0pt
\item the Muon momentum ramp, \texttt{min(step/300,1)}, completing at $15.0\%$ of the run
      instead of $30.3\%$;
\item the first eleven steps, which are uncharged, giving the control $5.8$M free tokens
      against the treatment's $2.9$M;
\item Adam bias correction (small: complete by step $50$--$300$);
\item Adam and NorMuon EMA horizons, fixed per step and so halved relative to the run;
\item the per-step weight-decay dose, doubled at the same wall clock;
\item gradient accumulation, $2\to1$.
\end{enumerate}

We first identified only the momentum ramp and queued two arms to decompose it. An
independent audit enumerated the rest, and those two arms bound one leg of six. Only leg
2 has a known direction, and it favours the \emph{control}: the measured win is if
anything understated.

We report this at length because the temptation was to publish $t=-60.7$ and describe a
critical-batch-size effect. The number is real and the mechanism is not established, and
at a fixed-time budget on this trainer, ``batch size'' and ``schedule shape'' may not be
separable by configuration alone.

\section{Threats to validity}

\textbf{Single seed.} Every result is seed 42 on one host, by operator direction to
conserve wall clock. Nothing here is established as seed-general.

\textbf{Degrees of freedom.} All $t$ statistics carry 3 df or fewer.

\textbf{Device advantage.} The headline best model ran on our fastest device, which
carries a $\sim$0.0025 advantage over the slowest. The \emph{effects} are
device-counterbalanced and unaffected; the headline \emph{number} is flattered.

\textbf{Instrument defects dominate.} Of 47 registered lessons, 8 are scientific negatives
and the remainder are process or measurement failures.

\section{What the instrument cost}
\label{sec:instrument}

We describe these because they are more transferable than the results, and because each
was found by audit rather than intuition.

\textbf{An instrument that could not measure what it claimed.} A GPU-slack metric built
from CUDA events bracketing the step. Such a pair measures elapsed time between two stream
markers \emph{including} idle intervals, and our bracket enclosed the data loader, pinning
the quantity to step time by construction. It read near-zero slack on every run, which
would have been reported as ``the GPU is saturated''.

\textbf{A metric silently discarded from every run.} The per-token FLOP estimate was
printed by every variant and dropped by a parser requiring whitespace after a colon where
the print emits none---absent from all 116 records. Recovered by re-reading each run's own
log and verified to reproduce from first principles.

\textbf{Three vacuous activation diagnostics.} In one session we shipped three hypotheses
whose diagnostic the control also satisfied, all registered \emph{after} the lesson
recording that failure mode existed. Prose did not prevent recurrence; a pre-registration
check testing a proposed diagnostic against the existing control corpus does.

\textbf{A policy that strangled the search.} The explore/exploit controller retired an axis
after four consecutive non-improving runs. A counterbalanced quad is four runs of the
\emph{same} value---so the moment we adopted the design that makes results trustworthy,
every axis began closing after one experiment. Three axes closed having tried exactly one
distinct value, including one that closed on evidence its \emph{opposite} direction was
worth testing.

\textbf{A ledger that could not reach its own threshold, then looked forward in time.} The
explore/exploit share was first measured over runs; since each experiment is one novel run
plus $\sim$7 counterbalancing replicates, the exploitation share could not exceed
$\sim$12.5\% against a 30\% floor. Corrected to per-experiment, it then judged whether an
axis ``pays'' using \emph{final-state} knowledge, retroactively classifying the runs that
\emph{discovered} each lever as exploitation of it, and reporting a balanced 50\% where the
true chronological share was 1/8. Three defects, in a ledger built specifically to keep us
honest about that number; each time the measurement flattered the thing it measured.

\textbf{A gate that could never be satisfied.} Adoption of the confirmed stack was deferred
pending a second-seed replication---a condition the operator's own single-seed directive
precludes. An audit named this a pocket veto and priced it: sixteen queued runs measured
against a control $0.004408$ worse than the demonstrated configuration. A gate that cannot
be satisfied is not caution; it is a refusal wearing caution's clothing, and it is worse
than an explicit refusal because it looks like a pending decision and is never revisited.

\paragraph{The pattern.} Every one of these fixes was discovered when a rule blocked
something we wanted to run. We have never once gone looking for a rule that was too
\emph{permissive}. That asymmetry is itself a threat to validity, and we do not know how to
correct it from the inside---which is the strongest argument we can make for independent,
adversarial audit as a standing component of autonomous research rather than an occasional
one.

\section{Conclusion}

Five interventions lower $\mathrm{val\_bpb}$ at a fixed 300-second budget, improving the
best model from $0.990345$ to $0.984017$; three of them combine to about $89\%$ of their
additive sum. The largest is reported as a configuration result whose mechanism we cannot
yet name. Three others do not, one by a very large margin. The measurement design making
these statements possible---concurrent yoked pairs, device counterbalancing,
pre-registered activation diagnostics, and a willingness to refuse a verdict when the
design breaks---mattered more than any individual intervention, because the effects are
roughly a twentieth of the noise an unpaired comparison would face.

\end{multicols}
\end{document}
